\Subhead{2.3.2}{Магнитоактивные среды. Тензор проницаемости плазмы в магнитном поле. Нормальные волны и их поляризация}

Магнитное поле связывает поперечные компоненты электронной скорости, поэтому круговые поляризации испытывают разные отклики. В бесстолкновительной холодной электронной плазме с неподвижными ионами
\begin{equation}
 -\ii\omega m\bm v=-e(\bm E+\bm v\times\bm B_0),\quad
 \omega_p^2=\frac{n_ee^2}{m\varepsilon_0},\quad \Omega_e=-\frac{eB_0}{m}.
\end{equation}
При $\bm B_0\parallel z$ решение линейной системы даёт
\begin{equation}
 \varepsilon=\begin{pmatrix}S&-\ii D&0\\\ii D&S&0\\0&0&P\end{pmatrix},\quad
 S=1-\frac{\omega_p^2}{\omega^2-\Omega_e^2},\quad
 D=\frac{\Omega_e\omega_p^2}{\omega(\omega^2-\Omega_e^2)},\quad
 P=1-\frac{\omega_p^2}{\omega^2}.
\end{equation}
Знак $D$ зависит от определения циклотронной частоты и соглашения о преобразовании Фурье; здесь $\Omega_e$ отрицательна. При $\bm k\parallel\bm B_0$ собственные векторы $(1,\pm\ii,0)$ имеют $n_\pm^2=S\pm D$. Разность фазовых набегов вращает плоскость линейной поляризации:
\begin{equation}
 \theta_F=\frac{\omega l}{2c}(n_+-n_-)
\end{equation}
при согласованном выборе знака угла. При $\bm k\perp\bm B_0$ обыкновенная мода с $\bm E\parallel\bm B_0$ имеет $n^2=P$, необыкновенная — $n^2=(S^2-D^2)/S$.

Резонансы и отсечки возникают из нулей и полюсов этих выражений. Столкновения делают тензор комплексным; тепловое движение добавляет пространственную дисперсию и кинетические эффекты. Поэтому холодная модель не описывает затухание Ландау.
